% ============================================================
%  conclusion.tex — TH-GNN paper, NeurIPS 2025
% ============================================================

\section{Conclusion}
\label{sec:conclusion}

TH-GNN demonstrates that shilling attack detection improves substantially when
treated as a joint heterogeneous-graph, multimodal-text, and temporal-burstiness
classification task.
The three signal streams—a two-layer HGT backbone with learnable sinusoidal
temporal encoding, cross-modal attention over review and item text, and a GRU
over log-inter-arrival times—address orthogonal failure modes of prior detectors:
the temporal stream catches coordinated injection campaigns that text-only methods
miss; the text streams expose cross-modal inconsistencies that graph-only methods
miss; and the heterogeneous graph encoding preserves type-specific fingerprints
that homogeneous GNNs discard.
Across five attack families and four datasets, TH-GNN achieved a grand-mean
F\textsubscript{1} of 0.870, with a $+$10.9\,pp margin over the strongest
text-only baseline on LLM-agent attacks—a margin that widened to $+$11.5\,pp at
the lowest injection rate, precisely where single-modality detectors are most
vulnerable.

Two limitations warrant attention.
First, the temporal stream requires per-event timestamps; platforms that expose
only aggregate interaction counts would reduce TH-GNN to a text-and-graph
detector (the ablation study estimates a $-$5.1\,pp cost for removing the
temporal stream entirely).
Second, LLM-generated profiles will evolve as models improve, and periodic
retraining on fresh adversarial samples will be necessary to sustain detection
accuracy against next-generation agents.
Future work should investigate adaptive attacks that specifically target temporal
coordination signals, online detection settings where profiles must be scored
incrementally as reviews arrive, and cross-platform transfer where a detector
trained on one review domain is applied to another.
